\section*{Chapitre \ref{ch:g11:lenses-and-eye} --- Lentilles, images et l'œil}

\begin{solution}{exo:g11:lenses-and-eye:1}
$C = 1/0.50 = +2.0$~D~; $1/(-0.20) = -5.0$~D~; $1/0.0040 = +250$~D.
Une \omterm{def:g11:lenses-and-eye:lens}{lentille} de $+8.0$~D~:
$f' = 1/8.0 = \qty{0.125}{m} = \qty{12.5}{cm}$.
\end{solution}

\begin{solution}{exo:g11:lenses-and-eye:2}
\omterm{def:g11:lenses-and-eye:lens}{Divergente} ($f' < 0$). Non~: pour toute distance
objet $d$, $\gamma = f'/(f' - d)$ (montré dans le
\cref{exo:g11:lenses-and-eye:15}) est entre $0$ et $1$ --- l'image est
toujours \omterm{def:g11:lenses-and-eye:images}{virtuelle}, droite et \emph{réduite}.
L'imprimé paraît droit et rétréci, jamais grossi.
\end{solution}

\begin{solution}{exo:g11:lenses-and-eye:3}
L'objet est en $2f'$~; les rayons central et parallèle se croisent à
\qty{20}{cm} derrière la \omterm{def:g11:lenses-and-eye:lens}{lentille}. Vérification~:
$1/\overline{OA'} = 1/10 - 1/20 = 1/20$, donc
$\overline{OA'} = +\qty{20}{cm}$ et $\gamma = 20/(-20) = -1$~:
\omterm{def:g11:lenses-and-eye:images}{réelle}, renversée, grandeur nature --- comme dessiné.
\end{solution}

\begin{solution}{exo:g11:lenses-and-eye:4}
$1/\overline{OA'} = 1/8.0 - 1/24 = 1/12$~:
$\overline{OA'} = +\qty{12}{cm}$, $\gamma = 12/(-24) = -0.50$.
\omterm{def:g11:lenses-and-eye:images}{Image réelle} à \qty{12}{cm} derrière la
\omterm{def:g11:lenses-and-eye:lens}{lentille}, renversée, moitié de taille.
\end{solution}

\begin{solution}{exo:g11:lenses-and-eye:5}
La \omterm{def:g11:lenses-and-eye:eye}{cornée} et le \omterm{def:g11:lenses-and-eye:eye}{cristallin}
forment la \omterm{def:g11:lenses-and-eye:lens}{lentille convergente}~; la
\omterm{def:g11:lenses-and-eye:eye}{rétine} est l'écran.
L'\omterm{def:g11:lenses-and-eye:eye}{accommodation} change la
\omterm{def:g11:lenses-and-eye:focal}{distance focale}
(\omterm{def:g11:lenses-and-eye:vergence}{vergence}) du
\omterm{def:g11:lenses-and-eye:eye}{cristallin}~; la distance cristallin--rétine ne peut
pas changer. \omterm{def:g11:lenses-and-eye:eye}{Punctum proximum}~: point le plus proche
vu nettement à pleine \omterm{def:g11:lenses-and-eye:eye}{accommodation} ---
\qty{25}{cm} pour l'\omterm{def:g11:lenses-and-eye:eye}{œil normal}.
\end{solution}

\begin{solution}{exo:g11:lenses-and-eye:6}
Paysage~: capteur dans le plan focal, à \qty{50}{mm} de la
\omterm{def:g11:lenses-and-eye:lens}{lentille}. À \qty{3.0}{m}~:
$1/\overline{OA'} = 1/50 - 1/3000 = 59/3000$, donc
$\overline{OA'} \approx \qty{50.8}{mm}$. Débattement~: environ
\qty{0.8}{mm}.
\end{solution}

\begin{solution}{exo:g11:lenses-and-eye:7}
$1/\overline{OA'} = 1/6.0 - 1/5.0 = -1/30$~:
$\overline{OA'} = \qty{-30}{cm}$ --- \omterm{def:g11:lenses-and-eye:images}{virtuelle},
droite, $\gamma = (-30)/(-5.0) = +6.0$~: six fois le timbre, à
\qty{30}{cm} au-dessus de la \omterm{def:g11:lenses-and-eye:lens}{lentille}.
\end{solution}

\begin{solution}{exo:g11:lenses-and-eye:8}
$1/f' = 1/40 + 1/40 = 1/20$~: $f' = \qty{20}{cm}$~;
$\gamma = 40/(-40) = -1$. La symétrie
\omterm{def:g10:inertia:force}{impose}
$\abs{\overline{OA'}} = \abs{\overline{OA}}$, d'où
$\abs{\gamma} = 1$~; l'image est \omterm{def:g11:lenses-and-eye:images}{réelle} et
renversée, donc $\gamma = -1$~: exactement grandeur nature.
\end{solution}

\begin{solution}{exo:g11:lenses-and-eye:9}
$1/\overline{OA'} = -1/25 - 1/50 = -3/50$~:
$\overline{OA'} \approx \qty{-16.7}{cm}$,
$\gamma = (-16.7)/(-50) = +0.33$~: \omterm{def:g11:lenses-and-eye:images}{virtuelle},
droite, un tiers de la taille. Les lunettes d'un myope sont
\omterm{def:g11:lenses-and-eye:lens}{divergentes}~: tout ce qu'on voit à travers --- y
compris, de l'extérieur, les yeux du porteur --- est réduit.
\end{solution}

\begin{solution}{exo:g11:lenses-and-eye:10}
Image de l'infini au \omterm{def:g11:lenses-and-eye:eye}{punctum remotum}~:
$f' = \qty{-0.40}{m}$, donc $C = -2.5$~D. L'image de la montagne se
trouve à \qty{40}{cm} devant la \omterm{def:g11:lenses-and-eye:lens}{lentille} ---
exactement le \omterm{def:g11:lenses-and-eye:eye}{punctum remotum}, que l'œil détendu voit
nettement avec zéro \omterm{def:g11:lenses-and-eye:eye}{accommodation}.
\end{solution}

\begin{solution}{exo:g11:lenses-and-eye:11}
$C = 1/\overline{OA'} - 1/\overline{OA} = 1/(-0.80) - 1/(-0.25)
= -1.25 + 4.00 = +2.75$~D.
\end{solution}

\begin{solution}{exo:g11:lenses-and-eye:12}
\emph{1.} $\gamma = \overline{OA'}/\overline{OA} = -50$ donne
$\overline{OA'} = -50\,\overline{OA}$~; avec
$\overline{OA'} - \overline{OA} = \qty{5.1}{m}$,
$-51\,\overline{OA} = \qty{5.1}{m}$~: $\overline{OA} = \qty{-0.10}{m}$,
$\overline{OA'} = +\qty{5.0}{m}$.

\emph{2.} $1/f' = 1/5.0 + 1/0.10 = 10.2$~: $f' \approx \qty{9.8}{cm}$.

\emph{3.} $50 \times \qty{24}{mm} = \qty{1.2}{m}$, renversée --- d'où
l'habitude de charger les diapositives à l'envers.
\end{solution}

\begin{solution}{exo:g11:lenses-and-eye:13}
\emph{1.} $C_\infty = 1/0.017 \approx 59$~D~;
$C_{25} = 1/0.017 + 1/0.25 \approx 63$~D.
\omterm{def:g10:signals-and-waves:amplitude}{Amplitude}~: $1/0.25 = 4.0$~D.

\emph{2.} \omterm{def:g10:signals-and-waves:amplitude}{Amplitude} $= 1/1.0 = 1.0$~D~: un quart
($25\,\%$) des $4.0$~D de l'\omterm{def:g11:lenses-and-eye:eye}{œil normal} reste.
\end{solution}

\begin{solution}{exo:g11:lenses-and-eye:14}
\emph{1.} $f' = 0.25/3 \approx \qty{0.083}{m} = \qty{8.3}{cm}$.

\emph{2.} $\overline{OA'} = \qty{-25}{cm}$~:
$1/\overline{OA} = -1/25 - 3/25 = -4/25$, d'où
$\overline{OA} = \qty{-6.25}{cm}$ et $\gamma = +4.0$ --- mieux que le
$\times 3$ annoncé, qui suppose l'image à l'infini~; la ramener au
\omterm{def:g11:lenses-and-eye:eye}{punctum proximum} gagne une
\omterm{def:g10:orders-of-magnitude:unit}{unité}.
\end{solution}

\begin{solution}{exo:g11:lenses-and-eye:15}
$1/\overline{OA'} = 1/f' - 1/d = (d - f')/(f'd)$, donc
$\overline{OA'} = f'd/(d - f')$ et
$\gamma = \overline{OA'}/(-d) = f'/(f' - d)$. Si $d > f'$~:
$\overline{OA'} > 0$ (\omterm{def:g11:lenses-and-eye:images}{réelle}) et $\gamma < 0$
(renversée). Si $0 < d < f'$~: $\overline{OA'} < 0$
(\omterm{def:g11:lenses-and-eye:images}{virtuelle}) et
$\gamma = f'/(f' - d) > 1$~: droite et agrandie --- la configuration de
\omterm{def:g11:lenses-and-eye:magnifier}{loupe}.
\end{solution}

\begin{solution}{pb:g11:lenses-and-eye:1}
\textbf{1.} $1/\overline{OA'} = 1/20 - 1/60 = 1/30$~:
$\overline{OA'} = +\qty{30}{cm}$, $\gamma = -0.50$ ---
\omterm{def:g11:lenses-and-eye:images}{réelle}, renversée, moitié de taille.

\textbf{2.} $\overline{OA'} = +\qty{60}{cm}$, $\gamma = -2$~: les
positions \qty{30}{cm} et \qty{60}{cm} ont échangé leurs rôles, les
\omterm{def:g11:lenses-and-eye:magnification}{grandissements} sont réciproques --- la lumière
inversée suit le même chemin.

\textbf{3.} $1/\overline{OA'} = 1/20 - 1/10 = -1/20$~:
$\overline{OA'} = \qty{-20}{cm}$, $\gamma = +2$ ---
\omterm{def:g11:lenses-and-eye:images}{virtuelle}, droite, doublée.

\textbf{4.} Depuis l'infini l'image commence dans le plan focal~; à
mesure que l'objet approche $F$ l'image
\omterm{def:g11:lenses-and-eye:images}{réelle}, renversée, s'éloigne vers l'infini et
grandit~; à l'intérieur de $f'$ elle devient
\omterm{def:g11:lenses-and-eye:images}{virtuelle}, droite, agrandie, devant la
\omterm{def:g11:lenses-and-eye:lens}{lentille}.

\textbf{5.} $1/\overline{OA} \approx 0$ laisse
$1/\overline{OA'} = 1/f'$~: $\overline{OA'} = f'$, le plan focal.

\textbf{6.} Dans le plan focal, à \qty{50}{mm} derrière la
\omterm{def:g11:lenses-and-eye:lens}{lentille}.

\textbf{7.} $1/\overline{OA'} = 1/50 - 1/1000 = 19/1000$~:
$\overline{OA'} \approx \qty{52.6}{mm}$~; débattement
$\approx \qty{2.6}{mm}$.

\textbf{8.} $1/\overline{OA} = 1/55 - 1/50 = -1/550$~: distance minimale
de mise au point \qty{550}{mm} = \qty{55}{cm}.

\textbf{9.} $\gamma = 55/(-550) = -0.10$~: le visage se forme à
\qty{20}{mm}, juste dans le capteur de \qty{24}{mm}.

\textbf{10.} $\gamma = -1$ \omterm{def:g10:inertia:force}{impose}
$\overline{OA'} = -\overline{OA}$, donc $2/\overline{OA'} = 1/f'$~:
objet à \qty{100}{mm} ($= 2f'$) devant, capteur à \qty{100}{mm}
derrière --- le double du tirage paysage, bien au-delà du mécanisme de
\qty{55}{mm}~; les objectifs macro fournissent le débattement
supplémentaire.

\textbf{11.} Tout ce qui est plus près que son
\omterm{def:g11:lenses-and-eye:eye}{punctum proximum} (\qty{1.0}{m}) est flou~: les bras
tendus amènent la page aussi près de \qty{1.0}{m} que possible ---
nette, si petite.

\textbf{12.} $C = 1/(-1.0) - 1/(-0.25) = -1.0 + 4.0 = +3.0$~D.

\textbf{13.} $\gamma = (-1.0)/(-0.25) = +4.0$~: quatre fois plus grande
mais quatre fois plus loin --- même taille angulaire~; le gain n'est
pas la taille mais la \emph{netteté}, puisque l'image se trouve
désormais à son \omterm{def:g11:lenses-and-eye:eye}{punctum proximum}.

\textbf{14.} La netteté exige l'image au-delà de \qty{1.0}{m}~: de la
page à \qty{25}{cm} (image à \qty{1.0}{m}) à la page en
$f' = 1/3.0 \approx \qty{0.33}{m}$ (image à l'infini). Un objet
lointain se forme à \qty{33}{cm} \emph{derrière} les lunettes --- une
\omterm{def:g11:lenses-and-eye:images}{image réelle} que l'œil ne peut utiliser~: le monde
lointain est flou.

\textbf{15.} D'où les bifocaux~: puissance de lecture
(\omterm{def:g10:energy-conservation:power}{puissance}) dans la moitié inférieure de la
\omterm{def:g11:lenses-and-eye:lens}{lentille}, aucune (ou la correction de loin) dans la
moitié supérieure.

\textbf{16.} $1/\overline{OA'} = 1/5.0 - 1/4.0 = -1/20$~:
$\overline{OA'} = \qty{-20}{cm}$, \omterm{def:g11:lenses-and-eye:images}{virtuelle},
droite, $\gamma = +5.0$.

\textbf{17.} $1/\overline{OA} = -1/25 - 1/5.0 = -6/25$~:
$\overline{OA} \approx \qty{-4.2}{cm}$~; $\gamma = (-25)/(-25/6) = +6.0$.

\textbf{18.} Pierre en $F$~: image à l'infini, vue avec zéro
\omterm{def:g11:lenses-and-eye:eye}{accommodation} --- des heures d'inspection sans fatigue
oculaire.

\textbf{19.} $C_\infty = 1/0.017 \approx 59$~D~;
$C_{25} = 1/0.017 + 1/0.25 \approx 63$~D~: environ $4$~D
d'\omterm{def:g11:lenses-and-eye:eye}{accommodation}.

\textbf{20.} $f' = 1/58.8 \approx \qty{17.0}{mm}$ contre
$1/62.8 \approx \qty{15.9}{mm}$~: environ \qty{1.1}{mm}, quelque
$6\,\%$ --- le zoom intégré que l'appareil photo imitait avec
\qty{2.6}{mm} de débattement et que grand-mère compensait par $+3$~D.
\end{solution}
